\subsection{Polysomnography and the Physiological Measurement Target}

Polysomnography (PSG) provides the reference framework for characterizing human
sleep because sleep stages are defined from multiple physiological signals rather
than from behavior or movement alone. Standard PSG combines
electroencephalography (EEG), electrooculography (EOG), and electromyography
(EMG), with additional respiratory and cardiovascular channels included when
required. Under conventional scoring rules, these signals are interpreted
together to classify successive 30-s epochs as wake, N1, N2, N3, or rapid eye
movement (REM) sleep \citep{silber2007scoring}.

For TMR, the relevance of PSG extends beyond the categorical stage label. EEG
also preserves information about transient electrophysiological events such as
slow waves, sleep spindles, K-complexes, and cortical arousals. These events can
occur at substantially shorter timescales than the epochs used for conventional
sleep staging. A recording may therefore be correctly classified as N2 or N3
while still containing moment-to-moment physiological variation relevant to
whether stimulation should occur.

PSG should nevertheless be treated as a standardized reference rather than as
an error-free representation of a continuously evolving brain state. Human sleep
scoring contains genuine ambiguity, particularly near stage transitions.
Rosenberg and Van Hout reported an average inter-scorer stage agreement of
82.6\% in the American Academy of Sleep Medicine inter-scorer program, with
lower agreement for N3 and especially N1 than for several other stages
\citep{rosenberg2013interscorer}. Wearable performance should consequently be
interpreted relative to a high-quality physiological reference whose own
categorical labels contain uncertainty.

The translational objective is therefore not to reproduce the complete PSG
laboratory in wearable form. It is to determine which components of the
physiological information available in PSG must remain observable for a
particular closed-loop decision. This distinction permits reduced sensing
systems to be evaluated according to functional adequacy rather than according
to whether they reproduce every channel used in clinical polysomnography.


\subsection{Wearable EEG: Preserving Electrophysiological Observability}

Wearable EEG attempts to preserve access to scalp electrophysiology while
reducing the recording burden of full PSG. Compared with peripheral sensing,
its principal scientific advantage is not that it measures sleep ``directly,''
but that the recorded signal contains electrophysiological features from which
sleep stages and several intervention-relevant events are conventionally
defined. Depending on montage and signal quality, wearable EEG may preserve
information about slow-wave activity, sleep spindles, stage transitions, and
cortical responses to sensory stimulation.

Reducing the number of electrodes introduces important trade-offs. Fewer
channels reduce spatial coverage and redundancy and may make it more difficult
to distinguish localized artifact from genuine physiology. They do not,
however, inherently remove EEG's high temporal resolution. The relevant
question is whether a particular montage preserves the signal features required
by the intended decision and whether those features remain reliable throughout
overnight use.

Empirical validation demonstrates both the potential and the limitations of
this approach. In a developer-associated validation of the Dreem headband,
Arnal et al.\ compared simultaneous headband and PSG recordings in 25
participants. Automatic sleep staging achieved a mean overall accuracy of
\(83.5\pm6.4\%\), compared with \(86.4\pm8.0\%\) average agreement for five
human sleep experts in the same study
\citep{arnal2020dreem}. These results showed that reduced-montage wearable EEG
can preserve substantial sleep-staging information under controlled recording
conditions.

Performance should not, however, be generalized from a single device,
population, or validation design. An independent evaluation of the same
headband in 62 older adults, including participants with Alzheimer's disease,
reported moderate five-stage epoch-level performance
(\(\mathrm{MCC}=0.54\pm0.14\)); after artifact removal, 73\% of recordings were
usable for quantitative EEG analysis
\citep{ravindran2025dreem}. The contrast between these studies illustrates that
wearable EEG performance depends on population, signal quality, recording
conditions, algorithm behavior, and the physiological quantity being extracted.

This evidence supports wearable EEG as a plausible means of retaining
electrophysiological observability outside full PSG, but not as an automatically
reliable solution. For closed-loop TMR, stage classification must therefore be
evaluated together with signal-quality detection, availability of continuous
data, preservation of intervention-relevant EEG features, and the ability to
withhold interpretation during corrupted intervals.


\subsection{PPG and Peripheral Sleep-State Inference}

Photoplethysmography (PPG) approaches sleep monitoring from a different
physiological level. PPG measures peripheral blood-volume changes associated
with the cardiac cycle, from which pulse rate, pulse-rate variability, waveform
characteristics, and respiratory-related information can be derived. Because
autonomic regulation changes across wakefulness, NREM sleep, and REM sleep,
these signals contain information that can be used to estimate sleep state
probabilistically \citep{dezambotti2019wearable}.

The distinction from EEG is therefore one of observability rather than simple
accuracy. PPG does not observe the cortical electrical phenomena used to define
slow waves, sleep spindles, or cortical arousals. Instead, an algorithm infers
sleep state from peripheral physiological correlates. Actigraphy, temperature,
respiration, electrodermal activity, and other peripheral channels can be added
to reduce ambiguity, but the resulting estimate remains an inference from
signals downstream from or correlated with the neural state being estimated.

This inferential distance does not imply that PPG-based sleep staging is
ineffective. Korkalainen et al.\ trained deep-learning models using PPG
recordings from 894 patients undergoing diagnostic PSG. In an independent test
set, three-stage classification of wake, NREM, and REM reached 80.1\% accuracy
with Cohen's \(\kappa=0.65\). Performance declined as the classification problem
became more specific: four-stage classification reached 68.5\%
(\(\kappa=0.54\)), while five-stage classification reached 64.1\%
(\(\kappa=0.51\))
\citep{korkalainen2020ppg}. These results demonstrate that peripheral cardiac
information contains substantial sleep-state information while also showing the
cost of resolving more specific physiological categories.

Consumer and multisensor devices show similarly variable results. Chinoy
et al.\ compared seven consumer sleep-tracking systems with simultaneous PSG in
healthy young adults. Sleep sensitivity was high across devices
(\(\geq0.93\)), but wake specificity ranged from 0.18 to 0.54, sleep-stage
performance was inconsistent, and several devices performed worse during
disrupted sleep
\citep{chinoy2021consumer}. A more recent multi-night study of the Oura Ring
Generation~3 reported strong agreement with PSG for several global sleep
measures and stage-specific classification, although the study received
financial support from the device manufacturer and should be interpreted in
that context \citep{svensson2024oura}. Together, these studies show that
peripheral sensing can provide useful and sometimes strong retrospective sleep
estimates, but performance is device-, algorithm-, population-, and
condition-dependent.

Importantly, peripheral sensing has already progressed beyond retrospective
sleep summaries in TMR research. SleepStim used smartwatch-derived heart-rate
and inertial signals to estimate the probability of N3 sleep in real time and
to govern auditory cue delivery at home \citep{whitmore2022improving}. This
constitutes direct proof of concept that peripheral physiological information
can participate in automated TMR control. It does not establish equivalence to
EEG-based control, nor does it demonstrate that general consumer sleep-stage
outputs are automatically suitable for TMR. Instead, it shows that the
scientific question must be framed at the level of the complete decision
process rather than sensor type alone.


\begin{table*}[t]

\centering

\caption{
Comparison of sleep-monitoring modalities relevant to closed-loop TMR.
The table distinguishes the physiological information observed by each modality
from the inference or intervention decision subsequently made from that
information.
}

\label{tab:sensing_modalities_conceptual}

\footnotesize
\setlength{\tabcolsep}{4pt}
\renewcommand{\arraystretch}{1.22}

\begin{tabularx}{\textwidth}{
@{}
>{\raggedright\arraybackslash}p{2.1cm}
>{\raggedright\arraybackslash}X
>{\raggedright\arraybackslash}X
>{\raggedright\arraybackslash}X
@{}
}

\toprule

\textbf{Modality}
&
\textbf{Physiological information}
&
\textbf{Representative evidence}
&
\textbf{TMR-relevant limitation}
\\

\midrule

\textbf{PSG}
&
Multimodal EEG, EOG, and EMG, with additional physiological channels as
required.
&
Reference framework for standardized sleep staging and observation of
stage-specific electrophysiology
\citep{silber2007scoring}.
&
High measurement burden; conventional stage labels remain epoch-based and
contain scorer uncertainty.
\\

\addlinespace[1.2mm]

\textbf{Wearable EEG}
&
Reduced-montage scalp electrophysiology with device-dependent access to
slow-wave, spindle, arousal, and stage-related features.
&
Dreem validation reported \(83.5\%\) mean staging accuracy in one controlled
study; performance was more moderate in an independent older population
\citep{arnal2020dreem,ravindran2025dreem}.
&
Reduced spatial redundancy and vulnerability to overnight artifact and
electrode-contact loss; performance is configuration- and population-specific.
\\

\addlinespace[1.2mm]

\textbf{PPG / cardiac}
&
Peripheral pulse and autonomic dynamics associated with sleep-state changes.
&
PPG-only models have demonstrated moderate-to-substantial agreement with PSG,
with performance declining as stage specificity increases
\citep{korkalainen2020ppg}.
&
Does not directly observe cortical slow waves, spindles, or cortical arousal;
stage must be inferred from peripheral correlates.
\\

\addlinespace[1.2mm]

\textbf{Multimodal peripheral}
&
PPG combined with movement and potentially temperature, respiration, or other
peripheral channels.
&
Consumer systems can provide useful sleep estimates, but performance varies
substantially across devices and recording conditions
\citep{chinoy2021consumer,svensson2024oura}.
&
Additional signals may improve stage inference without establishing
electrophysiological equivalence or intervention readiness.
\\

\addlinespace[1.2mm]

\textbf{Actigraphy}
&
Movement and sustained immobility measured through accelerometry.
&
Well suited to longitudinal sleep--wake estimation but provides limited
physiological stage information \citep{dezambotti2019wearable}.
&
Insufficient alone for distinguishing the electrophysiologically defined NREM
states required by stage-aware TMR.
\\

\bottomrule

\end{tabularx}

\vspace{1mm}

\begin{minipage}{\textwidth}
\footnotesize
\textit{Note.}
Performance values are not intrinsic properties of sensor modalities.
They depend on device configuration, preprocessing, classification algorithm,
population, recording environment, reference scoring procedure, and validation
design.
\end{minipage}

\end{table*}


\subsection{Validation Evidence and What Accuracy Means}

Validation of wearable sleep monitoring generally relies on simultaneous
recording with PSG followed by epoch-level or summary-level comparison. Such a
comparison is necessary but not sufficient for interpretation. The apparent
performance of a system depends on the definition of the classification task,
the composition of the validation population, the independence of the test
data, the quality of the reference scoring, and the metrics used to summarize
error \citep{dezambotti2019wearable}.

The number of predicted classes is particularly important. Sleep--wake
classification, three-stage classification, four-stage classification, and
full five-stage scoring are progressively different problems. Results obtained
for one should not be presented as evidence for another. The PPG results
reported by Korkalainen et al., for example, decreased from 80.1\% accuracy for
three-stage classification to 64.1\% for five-stage classification
\citep{korkalainen2020ppg}. A single headline accuracy value therefore has
little meaning unless the target classes are specified.

Overall accuracy can also conceal stage-specific failure. Sensitivity indicates
how often a true target state is recognized; precision indicates how often a
predicted target state is actually correct; specificity quantifies rejection of
non-target states; and confusion matrices reveal which physiological states are
systematically mistaken for one another. Cohen's kappa, Matthews correlation
coefficient, macro-F1, and related measures can reduce some distortions caused
by class imbalance, although no single metric captures every operational
consequence.

For TMR, false-negative and false-positive errors have different implications.
Failing to recognize one eligible NREM interval may simply remove an opportunity
for stimulation. Incorrectly declaring an unsuitable interval eligible may
cause a cue to be delivered outside the physiological conditions supported by
the evidence. The relative importance of sensitivity and precision must
therefore be determined from the intervention policy rather than from generic
sleep-staging conventions.

Generalization is equally important. Consumer-device studies show that
performance can deteriorate during disrupted sleep and vary across algorithms
and devices \citep{chinoy2021consumer}. Wearable EEG evidence likewise changes
across populations and signal-quality conditions
\citep{arnal2020dreem,ravindran2025dreem}. Validation results should therefore
be interpreted as properties of a particular sensor--algorithm--population
combination, not as universal properties of ``EEG'' or ``PPG.''

Independence of validation also matters. Developer- or manufacturer-supported
studies can provide valuable peer-reviewed evidence, but their provenance should
remain visible and their results should be considered together with independent
validation. The purpose of this distinction is not to discard industry-linked
evidence, but to avoid treating a single favorable validation as sufficient for
a modality-level conclusion.


\subsection{From Retrospective Scoring to Real-Time Closed-Loop Control}

Retrospective sleep scoring and real-time intervention control solve different
problems. An offline algorithm may use complete-night context, future epochs,
temporal smoothing, repeated preprocessing, and manual artifact review before
assigning the final hypnogram. A causal controller has access only to data
already acquired and must make a decision before the physiological opportunity
for stimulation has passed.

This distinction does not imply that all non-automated TMR protocols are
scientifically invalid. Foundational TMR experiments often relied on
experimenters who monitored sleep physiology and manually controlled
stimulation. The limitation arises when an autonomous system uses a fixed or
unsupervised schedule without verifying that the sleeper remains in an eligible
state. For such a system, physiological state must influence whether cue
delivery proceeds, is delayed, or is withheld.

Latency becomes meaningful within this causal setting. End-to-end delay includes
signal acquisition, window formation, preprocessing, inference, decision logic,
communication, and physical cue onset. A classifier may correctly describe the
data from which its estimate was produced while that estimate becomes less
relevant before the sound is delivered. Longer contextual windows can improve
classification stability but also make the underlying evidence older. The
appropriate latency threshold therefore cannot be assigned from computational
speed alone; it must be evaluated relative to the temporal stability of the
physiological condition being targeted.

SleepStim illustrates that peripheral sensing can meet at least some of these
requirements in practice. Its smartwatch transmitted heart-rate and inertial
data approximately once per second, and the system reported roughly 1~s latency
between acquisition and availability of its N3 probability estimate. Cueing was
suspended when signal transmission failed or the estimated probability of N3
fell below predefined thresholds
\citep{whitmore2022improving}. This is stronger evidence for closed-loop
readiness than retrospective stage agreement alone because the physiological
estimate actually governed intervention.

Wearable EEG has likewise been used within real-time home TMR. Salfi et al.\
used a wearable EEG headband to identify slow-wave activity and trigger
vocabulary-associated cues during home sleep
\citep{salfi2025promoting}. This demonstrates that electrophysiological
information can also be incorporated into an autonomous home intervention.
Neither study, however, establishes a universal architecture or a definitive
performance threshold for future systems.

Uncertainty must also be treated as part of control rather than hidden behind
a forced stage label. Signal degradation, state transitions, disagreement
between signals, or low classifier confidence may leave the system without
sufficient evidence to justify stimulation. An explicit abstention state is
therefore preferable to mandatory classification when the cost of an
inappropriate cue exceeds the cost of missing one opportunity.

Finally, the measurement task continues after stimulation. A cue may leave the
physiological state largely unchanged, or it may produce movement, cortical
activation, autonomic change, a transition to lighter sleep, or awakening.
Post-cue information can therefore determine whether subsequent stimulation
should continue, be delayed, or be suspended. Closed-loop TMR should be
understood as a recurrent measurement--decision--intervention process rather
than as a sleep-stage classifier connected directly to audio output.


\begin{table*}[t]

\centering

\caption{
Operational distinction between retrospective sleep scoring, causal
real-time state estimation, and closed-loop TMR.
}

\label{tab:closed_loop_requirements}

\footnotesize
\setlength{\tabcolsep}{5pt}
\renewcommand{\arraystretch}{1.22}

\begin{tabularx}{\textwidth}{
@{}
>{\raggedright\arraybackslash}p{2.7cm}
>{\raggedright\arraybackslash}X
>{\raggedright\arraybackslash}X
>{\raggedright\arraybackslash}X
@{}
}

\toprule

\textbf{Property}
&
\textbf{Retrospective scoring}
&
\textbf{Causal state estimation}
&
\textbf{Closed-loop TMR}
\\

\midrule

\textbf{Objective}
&
Reconstruct sleep architecture after recording.
&
Estimate the current physiological state as sleep unfolds.
&
Use physiological evidence to authorize and adapt intervention.
\\

\addlinespace[1.2mm]

\textbf{Temporal context}
&
Past and future epochs may be available.
&
Only previously acquired data are available.
&
Causal estimates are combined with cue history and subsequent physiological
responses.
\\

\addlinespace[1.2mm]

\textbf{Latency}
&
Generally not intervention-limiting.
&
Estimate must remain relevant to the current state.
&
Complete sensing-to-cue delay must preserve the validity of the intervention
decision.
\\

\addlinespace[1.2mm]

\textbf{Uncertainty}
&
Ambiguous epochs may be reconsidered retrospectively.
&
Confidence and signal quality must be available when the estimate is produced.
&
Insufficient evidence should result in delayed or withheld stimulation.
\\

\addlinespace[1.2mm]

\textbf{Signal loss}
&
Corrupted intervals can be removed or reviewed later.
&
Reduces or prevents current-state estimation.
&
Should produce abstention until usable physiological evidence returns.
\\

\addlinespace[1.2mm]

\textbf{Post-cue monitoring}
&
Not applicable.
&
Observation may continue without affecting an intervention.
&
Physiological responses influence whether subsequent cueing continues,
changes, or stops.
\\

\bottomrule

\end{tabularx}

\vspace{1mm}

\begin{minipage}{\textwidth}
\footnotesize
\textit{Note.}
Closed-loop readiness is defined here as a translational operating requirement,
not as a standardized clinical sleep-staging metric.
\end{minipage}

\end{table*}


\subsection{Evidence Boundary for Sensing-Pathway Selection}

The wearable literature supports three conclusions relevant to the translation
problem. First, reduced-montage EEG can preserve substantial
electrophysiological information and can support automated sleep staging outside
full PSG, although signal quality and performance vary across devices,
populations, and recording conditions
\citep{arnal2020dreem,ravindran2025dreem}. Second, PPG and multimodal peripheral
signals contain sufficient information to estimate broad sleep state and, in
some systems, more detailed sleep architecture
\citep{korkalainen2020ppg,chinoy2021consumer,svensson2024oura}. Third,
retrospective agreement with PSG is not equivalent to successful real-time
intervention control.

The relevant comparison for TMR is therefore not ``EEG accuracy versus PPG
accuracy'' in the abstract. Candidate pathways must instead be evaluated
against the functions established in the preceding sections: recognition of
physiologically supported cueing conditions, causal availability of the
measurement, signal-quality awareness, bounded sensing-to-intervention delay,
conservative handling of uncertainty, and detection of physiological
disturbance following stimulation.

Existing home TMR studies demonstrate that both peripheral and EEG-based
approaches can participate in automated cueing
\citep{whitmore2022improving,salfi2025promoting}. The available evidence does
not, within this section alone, establish equivalence between them or determine
which should serve as the first-generation implementation pathway. That
decision requires comparison of evidentiary continuity, physiological
observability, implementation uncertainty, and validation risk. Section~6
performs that translation using the requirements derived here.